\documentclass[12pt]{article}
\usepackage[T1]{fontenc}
\usepackage[utf8]{inputenc}
\usepackage{lmodern}
\usepackage{textcomp}
\usepackage{enumitem}
\usepackage{geometry}
\geometry{margin=1in}
\begin{document}
\begin{center}
Answer Key - Economics - K-12 Level Assessment 3 \
\small (Answers are ordered exactly as in the question paper.)
\end{center}

\section*{Multiple Choice}
\begin{enumerate}[label=\arabic*.]
\item \textbf{Answer:} D
\\ \textit{Options:} A) hurts creditors who do not anticipate it.; B) hurts creditors who anticipate it.; C) hurts debtors.; D) both A and D are correct.
\\ \textit{Note:} The correct answer is justified because both options A and D accurately describe different aspects of inflation, highlighting its causes and effects on the economy, thus providing a comprehensive understanding of the concept.
\item \textbf{Answer:} D
\\ \textit{Options:} A) the demand curve will shift to the left.; B) the demand curve will shift to the right.; C) the supply curve will shift to the left.; D) the supply curve will shift to the right.
\\ \textit{Note:} Subsidizing the production of halogen headlights lowers production costs for manufacturers, encouraging them to produce more, which shifts the supply curve to the right.
\item \textbf{Answer:} B
\\ \textit{Options:} A) frictionally unemployed.; B) cyclically unemployed.; C) seasonally unemployed.; D) structurally unemployed.
\\ \textit{Note:} Sue is considered cyclically unemployed because her job loss is directly tied to the downturn in the economy during a recession, which reduces overall demand for goods and services, leading to layoffs.
\item \textbf{Answer:} C
\\ \textit{Options:} A) marginal cost (MC) intersects AVC and ATC at their maximum points.; B) AVC and ATC intersect MC at its maximum point.; C) MC intersects AVC and ATC at their minimum points.; D) AVC and ATC intersect MC at its minimum point.
\\ \textit{Note:} The marginal cost (MC) curve intersects the average variable cost (AVC) and average total cost (ATC) curves at their minimum points because, at these points, the cost of producing one more unit is exactly equal to the average cost of production, indicating that any additional production will not decrease the average costs further.
\item \textbf{Answer:} C
\\ \textit{Options:} A) increases in taxes to fight recessions.; B) decreases in taxes to fight inflations.; C) changes in government spending and taxes to fight recessions or inflations.; D) federal deficits.
\\ \textit{Note:} Fiscal policy involves the use of government spending and taxation as tools to influence the economy, particularly in addressing issues such as recessions and inflation.
\end{enumerate}

\section*{True / False}
\begin{enumerate}[label=\arabic*.]
\setcounter{enumi}{5}
\item \textbf{Answer:} False
\\ \textit{Options:} A) True; B) False
\\ \textit{Note:} The appropriate fiscal policy to remedy a recession typically involves the government increasing spending or cutting taxes, which often results in a budget deficit rather than a surplus to stimulate economic growth.
\item \textbf{Answer:} False
\\ \textit{Options:} A) True; B) False
\\ \textit{Note:} A competitive firm making zero economic profit is covering all its opportunity costs, which means it is earning normal profits, just like other firms in the market.
\item \textbf{Answer:} False
\\ \textit{Options:} A) True; B) False
\\ \textit{Note:} Decreasing taxes while keeping spending unchanged can increase disposable income and consumer spending, which may exacerbate an inflationary gap rather than effectively address it.
\item \textbf{Answer:} False
\\ \textit{Options:} A) True; B) False
\\ \textit{Note:} Low barriers to entry typically lead to more competitors in a market, which increases consumer choices and can enhance price elasticity of demand, making the statement false.
\item \textbf{Answer:} True
\\ \textit{Options:} A) True; B) False
\\ \textit{Note:} The statement is true because Marginal Social Costs account for both the costs incurred by the producer (Private Marginal Costs) and the additional costs imposed on society due to negative externalities, reflecting the true cost of production to society as a whole.
\end{enumerate}

\section*{Long Form Response}
\begin{enumerate}[label=\arabic*.]
\setcounter{enumi}{10}
\item \textbf{Answer:} Comparative advantage plays a significant role in international trade, as it determines which countries are more efficient in producing certain goods. In this case, the United States has a comparative advantage in producing salmon due to its favorable climate, advanced fishing technology, and abundant resources, allowing it to produce salmon more efficiently than many other countries. Conversely, Peru has a comparative advantage in anchovy production, leveraging its coastal access and favorable fishing conditions to harvest anchovies at a lower opportunity cost. By trading salmon for anchovies, both countries can specialize in what they produce best, leading to increased overall efficiency and greater availability of both products in the global market. This exchange ultimately benefits consumers in both nations by providing a wider variety of seafood options.
\\ \textit{Note:} correct because it illustrates how the United States' efficiency in producing salmon and Peru's efficiency in producing anchovies, due to their respective resources and conditions, exemplify the principle of comparative advantage, which drives international trade by enabling countries to specialize in and exchange goods they can produce more efficiently.
\item \textbf{Answer:} The establishment of a new software company by a Mexican entrepreneur in St. Louis can significantly influence aggregate demand in the United States. As the company grows, it will create jobs, leading to increased income for workers, which can boost consumer spending. This rise in spending stimulates demand for goods and services in the local economy, contributing to overall economic growth. Additionally, the company may attract investors and other businesses to the area, further enhancing economic activity. Overall, the entrepreneurial venture not only benefits St. Louis but can have a ripple effect on the national economy by increasing aggregate demand.
\\ \textit{Note:} correct because the establishment of the software company leads to job creation and increased income, which boosts consumer spending and stimulates aggregate demand, thereby contributing to economic growth in the United States.
\end{enumerate}

\vfill
\noindent\textit{End of Answer Key}
\end{document}